\section{VIEW: Membuat dan Menggunakan}

\textbf{VIEW} adalah query SELECT yang disimpan sebagai tabel virtual di database. View tidak menyimpan data sendiri --- setiap kali diakses, MySQL menjalankan query yang mendasarinya. View berguna untuk: (1) menyederhanakan query kompleks (terutama yang menggunakan JOIN); (2) membatasi akses data --- berikan user akses ke view, bukan ke tabel asli; (3) menyediakan interface yang konsisten meskipun struktur tabel berubah \cite{mysqlDocs}.

\begin{webcode}[language=SQL, caption={Membuat dan menggunakan VIEW}]
-- Membuat VIEW: daftar produk dengan nama kategori
CREATE VIEW v_produk_kategori AS
SELECT p.id, p.nama, p.harga, p.stok, k.nama AS kategori
FROM produk p
LEFT JOIN kategori k ON p.kategori_id = k.id;

-- Menggunakan view seperti tabel biasa
SELECT * FROM v_produk_kategori WHERE kategori = 'Elektronik';

-- VIEW untuk ringkasan
CREATE VIEW v_ringkasan_kategori AS
SELECT k.nama AS kategori, COUNT(p.id) AS jumlah, AVG(p.harga) AS rata_harga
FROM kategori k LEFT JOIN produk p ON k.id = p.kategori_id
GROUP BY k.nama;

-- Menghapus view
DROP VIEW IF EXISTS v_produk_kategori;
\end{webcode}

